\documentclass{article}

% This is the preamble, load any packages you're going to use here
\usepackage{physics} % provides lots of nice features and commands often used in physics, it also loads some other packages (like AMSmath)
%\usepackage{siunitx} % typesets numbers with units very nicely
\usepackage{rerunfilecheck}
\usepackage{enumerate} % allows us to customize our lists
\usepackage{amsmath} % provides many useful commands for math
\usepackage{amssymb} % provides many useful symbols
\usepackage{graphicx} % allows us to include graphics
\usepackage{hyperref} % allows us to include hyperlinks
\usepackage{float} % allows us to control the placement of figures
\usepackage{listings} % allows us to include code
\usepackage{color} % allows us to define colors
\usepackage{tikz} % allows us to draw pictures
\usepackage{pgfplots}
\pgfplotsset{compat=1.18} % 设置 pgfplots 兼容版本
\usepackage{listings}
\usepackage{graphicx}
\usepackage[a4paper, left=2cm, right=2cm, top=1cm, bottom=2cm]{geometry}% 设置页边距
\usepackage{ctex}
%\setCJKmainfont[ItalicFont=Noto Sans CJK SC Bold, BoldFont=Noto Serif CJK SC Black]{Noto Serif CJK SC}
\lstset{
    columns=fixed,          
    numberstyle=\tiny\color{gray},                       % 设定行号格式
    frame=none,                                          % 不显示背景边框
    backgroundcolor=\color[RGB]{245,245,244},            % 设定背景颜色
    keywordstyle=\color[RGB]{40,40,255},                 % 设定关键字颜色
    numberstyle=\footnotesize\color{darkgray},           
    commentstyle=\it\color[RGB]{0,96,96},                % 设置代码注释的格式
    stringstyle=\rmfamily\slshape\color[RGB]{128,0,0},   % 设置字符串格式
    showstringspaces=false,                              % 不显示字符串中的空格
    language=python,                                        % 设置语言
    basicstyle=\ttfamily\footnotesize,                          % 设置代码字体大小
}


\begin{document}
\title{极化源相关论文学习笔记}
\author{刘世珺}
\date{\today}
\maketitle
\hypersetup{hidelinks}
\tableofcontents



\newpage
\section{极化H/D源}  




\section{极化氢分子源}

\subsection{Stimulated Raman adiabatic passage in physics,  chemistry, and beyond\cite{vitanov_stimulated_2017}}
Stimulated Raman adiabatic passage (STIRAP) 是一种利用中间态，将两个辐射场与两个量子态耦合，进一步可以使两个量子态之间相互转化的技术。这种技术有两个特点：
\begin{enumerate}
    \item 不受中间态自发辐射的损失
    \item 在不同实验条件下都比较稳定
\end{enumerate} 
下面就不同部分对STIRAP进行介绍。
\paragraph{基本原理}

在量子力学中，含时薛定谔方程为 
\begin{equation}
    i\hbar\frac{\text{d}}{\text{d} t}\Psi(t) = \hat{H}\Psi(t)
\end{equation}  

\subsection{Spin manipulation and nuclear polarization enhancement in particle beams with static magnetic fields\cite{kannis_spin_2025}}

当非相对论性粒子束在通过一个静态空间变化的磁场时，磁场可以增强氢/氘原子束或者特定转动状态的氢氘分子束中的核极化。这种增强源于超精细结构的跃迁。 通过研究不同初始条件下，非均匀磁场下的相互作用角动量动力学，可以了解这种增强。  该论文主要介绍自选动力学的理论框架，并给出在特定原子和分子中的应用。

\paragraph{基本原理}
本文考虑的系统包含以下几种角动量组成：核自旋角动量\(\textbf{I}\),电子的总自旋角动量\(\textbf{S}\)，分子的旋转角动量\(\textbf{J}\).假设由两个相互作用的角动量\(\textbf{L}_{\textbf{1,2}}\)组成的系统未耦合的基向量可以写作
\begin{equation}
    \ket{L_1l_1,L_2l_2} = \ket{L_1l_1}\otimes\ket{L_2l_2}
    \label{eq:uncoupled_basis_vector}   
\end{equation}  
其中\(L_{1,2}\)分别表示两个角动量的量子数，\(l_{1,2}\)表示它们沿着量子化轴的投影。所以可以将\(\ket{l_1,l_2}\)记作系统未耦合的基向量。

对于耦合基，定义一个额外的角动量算符\( \textbf{K} = \textbf{L}_{\textbf{2}} + \textbf{L}_{\textbf{2}} \),这个可以表示为\( \textbf{K} = \textbf{L}_{\textbf{1}}\otimes \mathbb{I}  + \mathbb{I} \otimes\textbf{L}_{\textbf{2}} \)，类似的。\(\ket{K,k}\)表示总角动量量子数为\(K\)的基向量在量子化轴上的投影为\(k\)。 



\subsection{Highly Spin-Polarized Atoms and Molecules from Rotationally State-Selected Molecules\cite{rakitzis_highly_2005}}

产生极化原子的方法目前有三种,分别是Stern-Gerlach分离法,光泵浦法(optical pumping)和自旋交换光泵浦法(spin-exchange optical pumping)。本篇论文展示了如何使用光学方法产生100\%的核极化，使用脉冲的激光或者微波可以产生具有量子数$J$和$M$的旋转极化的分子，之后将旋转的极化耦合到非极化的核自旋上，通过超精细作用实现极化的传递。





\section{极化$^{3}\text{He}$源}

\newpage
\addcontentsline{toc}{section}{参考文献}
\bibliographystyle{plain}
\bibliography{mylib}

\end{document}
